One of the applications of the transformation theorem is that we can now define Sobolev spaces on manifolds. Let $M$ be a compact $p$-dimensional surface in $\mathbb{R}^m$. We assume that $M$ is smooth of class $C^k$, $k\geq1$. Then every point in $M$ has a neighborhood within which $M$ can be represented in the form $\bm{x}=\bm{g}(\bm{y})$, where $\bm{y}\in\mathbb{R}^p$ and $\bm{g}$ is of class $C^k$. We can cover $M$ with a finite number of such neighborhoods; let us call them $U_i$, $1\leq i\leq N$, and let $\bm{g}_i:N_i\to U_i\cap M$ be the corresponding parametric representations of the surface. We make the following definition:

\begin{Definition}
\label{def:closed-set-partition-of-unity}
Let $S$ be a closed subset of $\mathbb{R}^m$ and let $\{U_j\}$ be a locally finite covering of $S$ with bounded open subsets of $\mathbb{R}^m$ (not of $S$). A \textbf{partition of unity} subordinate to the covering $\{U_j\}$ is a set of functions $\phi_j\in\mathcal{D}(\mathbb{R}^m)$ such that\begin{enumerate}
\item $0\leq\phi_j\leq1$,
\item $\operatorname{supp}\phi_j\subset U_j$, and
\item $\sum_j\phi_j=1$ in a neighborhood of $S$.
\end{enumerate}
\end{Definition}

The proof of the existence of a partition of unity is analogous to that of Theorem~\ref{thm:partition-of-unity-existence}. Note that if $S$ is compact, we can reduce every covering $\{U_j\}$ to a finite one.

Let $\{\phi_i\}$ be a partition of unity on $M$ subordinate to the covering $\{U_i\}$.

\begin{Definition}
\label{def:manifold-sobolev-spaces}
Let $M$ be as described above and let $l\leq k$ be a non-negative integer. We say that $u\in H^l(M)$ if $\phi_i(u\circ\bm{g}_i)\in H^l(\mathbb{R}^p)$ for every $i=1,\ldots,N$.
\end{Definition}

Here, it is of course understood that $\phi_i(u\circ\bm{g}_i)$ is set equal to zero outside $N_i$. The transformation theorem guarantees that this definition depends only on the surface and not on the parametrizations chosen. An inner product on $H^l(M)$ can be defined naturally by
\begin{equation}\label{eq:manifold-sobolev-local-inner-product}
(u,v)_{M,k}=\sum_{i=1}^N(\phi_i(u\circ\bm{g}_i),\phi_i(v\circ\bm{g}_i))_k.
\end{equation}
This definition does of course depend on the parametrizations, but all the inner products defined in this way are equivalent (Problem~\ref{ex:manifold-sobolev-inner-products-equivalence}). We note that the definition of $H^l(M)$ is easily extended to fractional $l$, except that for this case we have not established a transformation theorem. We state such a theorem without proof.

\begin{Theorem}
\label{thm:fractional-order-pullback-transformation}
Let $0\leq l\leq k$ with $k\in\mathbb{N}$ and $l$ real. Let $\bm{\Phi}$ be a $k$-diffeomorphism of $\mathbb{R}^m$. Then the pullback operator $\bm{\Phi}^*$ is a bounded linear mapping from $H^l(\mathbb{R}^m)$ onto itself.
\end{Theorem}

A proof can be based either on the theory of interpolation spaces or on \eqref{eq:fractional-sobolev-intrinsic-inner-product}. In applications, we are interested in Sobolev spaces on manifolds primarily for boundary data, and $l$ is a half-integer. For $l$ a half-integer, we shall obtain Theorem~\ref{thm:fractional-order-pullback-transformation} later as a corollary of the trace theorem.
